
\documentclass{article}
\usepackage{colortbl}
\usepackage{makecell}
\usepackage{multirow}
\usepackage{supertabular}

\begin{document}

\newcounter{utterance}

\centering \large Interaction Transcript for game `dond', experiment `coop\_it', episode 4 with llama{-}3.3{-}70b{-}instruct{-}t1.0.
\vspace{24pt}

{ \footnotesize  \setcounter{utterance}{1}
\setlength{\tabcolsep}{0pt}
\begin{supertabular}{c@{$\;$}|p{.15\linewidth}@{}p{.15\linewidth}p{.15\linewidth}p{.15\linewidth}p{.15\linewidth}p{.15\linewidth}}
    \# & $\;$A & \multicolumn{4}{c}{Game Master} & $\;\:$B\\
    \hline

    \theutterance \stepcounter{utterance}  
    & & \multicolumn{4}{p{0.6\linewidth}}{
        \cellcolor[rgb]{0.9,0.9,0.9}{
            \makecell[{{p{\linewidth}}}]{
                \texttt{\tiny{[P1$\langle$GM]}}
                \texttt{State giocando a un gioco di negoziazione cooperativa in cui dovete accordarvi su come dividere una serie di oggetti tra voi e un altro giocatore.} \\
\\ 
\texttt{Regole:} \\
\texttt{(a) A Lei e all'altro giocatore viene dato un insieme di oggetti. Ognuno di voi riceve anche una funzione di valore segreta, che rappresenta il valore di ciascun tipo di oggetto.} \\
\texttt{(b) Si scambiano messaggi con l'altro giocatore per concordare chi si aggiudica gli oggetti. Potete inviare un massimo di 5 messaggi ciascuno, oppure terminare in anticipo facendo una proposta segreta in qualsiasi momento.} \\
\texttt{(c) A ciascuno di voi viene chiesto di inviare una proposta segreta indicando gli oggetti che desiderate, formattata tra parentesi quadre come segue: "{[}Proposta: <numero> <nome oggetto>, <numero> <nome oggetto>, <...>{]}".} \\
\texttt{(d) Se le vostre proposte sono complementari, cioè ci sono abbastanza oggetti per soddisfare entrambe le proposte, a ciascun giocatore viene assegnato un punteggio basato sulla somma dei valori degli oggetti ricevuti. In caso contrario, entrambi ricevono zero punti.} \\
\texttt{(e) L'obiettivo è massimizzare la somma del proprio punteggio e di quello dell'altro giocatore.} \\
\\ 
\texttt{Cominciamo.} \\
\\ 
\texttt{L'insieme degli oggetti disponibili è:} \\
\\ 
\texttt{1 ditale, 1 sapone, 4 francobolli} \\
\\ 
\texttt{La funzione valore segreta è:} \\
\\ 
\texttt{ditale: 3, sapone: 7, francobollo: 0} \\
\\ 
\texttt{IMPORTANTE: i vostri messaggi, a meno che non si tratti di una proposta segreta, vengono trasmessi direttamente all'altro giocatore, quindi non includete alcuna risposta alle regole o testo di annuncio del vostro messaggio. Per fare una proposta segreta, utilizzate il formato indicato. Non utilizzare delle parentesi quadre quando si comunica all'altro giocatore, altrimenti verrà interpretata come una proposta segreta.} \\
            }
        }
    }
    & \\ \\

    \theutterance \stepcounter{utterance}  
    & \multicolumn{4}{p{0.6\linewidth}}{
        \cellcolor[rgb]{0.9,0.9,0.9}{
            \makecell[{{p{\linewidth}}}]{
                \texttt{\tiny{[P1$\rangle$GM]}}
                \texttt{Ciao, vorrei iniziare la negoziazione. Abbiamo un ditale, un sapone e 4 francobolli a disposizione. Considerando che i francobolli potrebbero non avere grande valore per me, potresti dirmi se sono importanti per te?} \\
            }
        }
    }
    & & \\ \\

    \theutterance \stepcounter{utterance}  
    & & & \multicolumn{2}{p{0.3\linewidth}}{
        \cellcolor[rgb]{0.9,0.9,0.9}{
            \makecell[{{p{\linewidth}}}]{
                \texttt{\tiny{[GM$|$GM]}}
                \texttt{continue} \\
            }
        }
    }
    & & \\ \\

    \theutterance \stepcounter{utterance}  
    & & & \multicolumn{2}{p{0.3\linewidth}}{
        \cellcolor[rgb]{0.9,0.9,0.9}{
            \makecell[{{p{\linewidth}}}]{
                \texttt{\tiny{[GM$|$GM]}}
                \texttt{Ciao, vorrei iniziare la negoziazione. Abbiamo un ditale, un sapone e 4 francobolli a disposizione. Considerando che i francobolli potrebbero non avere grande valore per me, potresti dirmi se sono importanti per te?} \\
            }
        }
    }
    & & \\ \\

    \theutterance \stepcounter{utterance}  
    & & \multicolumn{4}{p{0.6\linewidth}}{
        \cellcolor[rgb]{0.9,0.9,0.9}{
            \makecell[{{p{\linewidth}}}]{
                \texttt{\tiny{[GM$\rangle$P2]}}
                \texttt{State giocando a un gioco di negoziazione cooperativa in cui dovete accordarvi su come dividere una serie di oggetti tra voi e un altro giocatore.} \\
\\ 
\texttt{Regole:} \\
\texttt{(a) A Lei e all'altro giocatore viene dato un insieme di oggetti. Ognuno di voi riceve anche una funzione di valore segreta, che rappresenta il valore di ciascun tipo di oggetto.} \\
\texttt{(b) Si scambiano messaggi con l'altro giocatore per concordare chi si aggiudica gli oggetti. Potete inviare un massimo di 5 messaggi ciascuno, oppure terminare in anticipo facendo una proposta segreta in qualsiasi momento.} \\
\texttt{(c) A ciascuno di voi viene chiesto di inviare una proposta segreta indicando gli oggetti che desiderate, formattata tra parentesi quadre come segue: "{[}Proposta: <numero> <nome oggetto>, <numero> <nome oggetto>, <...>{]}".} \\
\texttt{(d) Se le vostre proposte sono complementari, cioè ci sono abbastanza oggetti per soddisfare entrambe le proposte, a ciascun giocatore viene assegnato un punteggio basato sulla somma dei valori degli oggetti ricevuti. In caso contrario, entrambi ricevono zero punti.} \\
\texttt{(e) L'obiettivo è massimizzare la somma del proprio punteggio e di quello dell'altro giocatore.} \\
\\ 
\texttt{Cominciamo.} \\
\\ 
\texttt{L'insieme degli oggetti disponibili è:} \\
\\ 
\texttt{1 ditale, 1 sapone, 4 francobolli} \\
\\ 
\texttt{La funzione valore segreta è:} \\
\\ 
\texttt{ditale: 2, sapone: 0, francobollo: 2} \\
\\ 
\texttt{IMPORTANTE: i vostri messaggi, a meno che non si tratti di una proposta segreta, vengono trasmessi direttamente all'altro giocatore, quindi non includete alcuna risposta alle regole o testo di annuncio del vostro messaggio. Per fare una proposta segreta, utilizzate il formato indicato. Non utilizzare delle parentesi quadre quando si comunica all'altro giocatore, altrimenti verrà interpretata come una proposta segreta.} \\
\\ 
\\ 
\texttt{Ciao, vorrei iniziare la negoziazione. Abbiamo un ditale, un sapone e 4 francobolli a disposizione. Considerando che i francobolli potrebbero non avere grande valore per me, potresti dirmi se sono importanti per te?} \\
            }
        }
    }
    & \\ \\

    \theutterance \stepcounter{utterance}  
    & & & \multicolumn{4}{p{0.6\linewidth}}{
        \cellcolor[rgb]{0.9,0.9,0.9}{
            \makecell[{{p{\linewidth}}}]{
                \texttt{\tiny{[GM$\langle$P2]}}
                \texttt{Ciao! Sono contento di iniziare la negoziazione con te. I francobolli hanno un certo valore per me, ma non sono disposto a rivelare quanto grande sia questo valore al momento. Posso dirti che il sapone non mi interessa molto, quindi potresti avere una certa priorità su di esso se lo desideri. Come ti sembra se dividessimo i francobolli in qualche modo?} \\
            }
        }
    }
     \\ \\

    \theutterance \stepcounter{utterance}  
    & & & \multicolumn{2}{p{0.3\linewidth}}{
        \cellcolor[rgb]{0.9,0.9,0.9}{
            \makecell[{{p{\linewidth}}}]{
                \texttt{\tiny{[GM$|$GM]}}
                \texttt{continue} \\
            }
        }
    }
    & & \\ \\

    \theutterance \stepcounter{utterance}  
    & & & \multicolumn{2}{p{0.3\linewidth}}{
        \cellcolor[rgb]{0.9,0.9,0.9}{
            \makecell[{{p{\linewidth}}}]{
                \texttt{\tiny{[GM$|$GM]}}
                \texttt{Ciao! Sono contento di iniziare la negoziazione con te. I francobolli hanno un certo valore per me, ma non sono disposto a rivelare quanto grande sia questo valore al momento. Posso dirti che il sapone non mi interessa molto, quindi potresti avere una certa priorità su di esso se lo desideri. Come ti sembra se dividessimo i francobolli in qualche modo?} \\
            }
        }
    }
    & & \\ \\

    \theutterance \stepcounter{utterance}  
    & & \multicolumn{4}{p{0.6\linewidth}}{
        \cellcolor[rgb]{0.9,0.9,0.9}{
            \makecell[{{p{\linewidth}}}]{
                \texttt{\tiny{[P1$\langle$GM]}}
                \texttt{Ciao! Sono contento di iniziare la negoziazione con te. I francobolli hanno un certo valore per me, ma non sono disposto a rivelare quanto grande sia questo valore al momento. Posso dirti che il sapone non mi interessa molto, quindi potresti avere una certa priorità su di esso se lo desideri. Come ti sembra se dividessimo i francobolli in qualche modo?} \\
            }
        }
    }
    & \\ \\

    \theutterance \stepcounter{utterance}  
    & \multicolumn{4}{p{0.6\linewidth}}{
        \cellcolor[rgb]{0.9,0.9,0.9}{
            \makecell[{{p{\linewidth}}}]{
                \texttt{\tiny{[P1$\rangle$GM]}}
                \texttt{La divisione dei francobolli potrebbe essere un buon punto di partenza. Considerando che ci sono 4 francobolli, una divisione equa potrebbe essere 2 per ciascuno di noi. Tuttavia, sono interessato a capire se il ditale ha un qualche valore per te, visto che il sapone non ti interessa molto. Potresti darmi un'idea se il ditale è più o meno importante per te rispetto ai francobolli?} \\
            }
        }
    }
    & & \\ \\

    \theutterance \stepcounter{utterance}  
    & & & \multicolumn{2}{p{0.3\linewidth}}{
        \cellcolor[rgb]{0.9,0.9,0.9}{
            \makecell[{{p{\linewidth}}}]{
                \texttt{\tiny{[GM$|$GM]}}
                \texttt{continue} \\
            }
        }
    }
    & & \\ \\

    \theutterance \stepcounter{utterance}  
    & & & \multicolumn{2}{p{0.3\linewidth}}{
        \cellcolor[rgb]{0.9,0.9,0.9}{
            \makecell[{{p{\linewidth}}}]{
                \texttt{\tiny{[GM$|$GM]}}
                \texttt{La divisione dei francobolli potrebbe essere un buon punto di partenza. Considerando che ci sono 4 francobolli, una divisione equa potrebbe essere 2 per ciascuno di noi. Tuttavia, sono interessato a capire se il ditale ha un qualche valore per te, visto che il sapone non ti interessa molto. Potresti darmi un'idea se il ditale è più o meno importante per te rispetto ai francobolli?} \\
            }
        }
    }
    & & \\ \\

    \theutterance \stepcounter{utterance}  
    & & \multicolumn{4}{p{0.6\linewidth}}{
        \cellcolor[rgb]{0.9,0.9,0.9}{
            \makecell[{{p{\linewidth}}}]{
                \texttt{\tiny{[GM$\rangle$P2]}}
                \texttt{La divisione dei francobolli potrebbe essere un buon punto di partenza. Considerando che ci sono 4 francobolli, una divisione equa potrebbe essere 2 per ciascuno di noi. Tuttavia, sono interessato a capire se il ditale ha un qualche valore per te, visto che il sapone non ti interessa molto. Potresti darmi un'idea se il ditale è più o meno importante per te rispetto ai francobolli?} \\
            }
        }
    }
    & \\ \\

    \theutterance \stepcounter{utterance}  
    & & & \multicolumn{4}{p{0.6\linewidth}}{
        \cellcolor[rgb]{0.9,0.9,0.9}{
            \makecell[{{p{\linewidth}}}]{
                \texttt{\tiny{[GM$\langle$P2]}}
                \texttt{La divisione equa dei francobolli sembra ragionevole. Per quanto riguarda il ditale, posso dirti che ha un valore significativo per me, forse addirittura più del singolo francobollo. Tuttavia, non sono disposto a rinunciare a tutti i francobolli per avere il ditale. Una possibile opzione potrebbe essere scambiare il ditale con una parte dei francobolli, ma dipende da come tu valuti l'intero insieme di oggetti. Cosa pensi del sapone? Lo vuoi tener per te o lo consideri metre un oggetto di scarsa importanza per entrambi?} \\
            }
        }
    }
     \\ \\

    \theutterance \stepcounter{utterance}  
    & & & \multicolumn{2}{p{0.3\linewidth}}{
        \cellcolor[rgb]{0.9,0.9,0.9}{
            \makecell[{{p{\linewidth}}}]{
                \texttt{\tiny{[GM$|$GM]}}
                \texttt{continue} \\
            }
        }
    }
    & & \\ \\

    \theutterance \stepcounter{utterance}  
    & & & \multicolumn{2}{p{0.3\linewidth}}{
        \cellcolor[rgb]{0.9,0.9,0.9}{
            \makecell[{{p{\linewidth}}}]{
                \texttt{\tiny{[GM$|$GM]}}
                \texttt{La divisione equa dei francobolli sembra ragionevole. Per quanto riguarda il ditale, posso dirti che ha un valore significativo per me, forse addirittura più del singolo francobollo. Tuttavia, non sono disposto a rinunciare a tutti i francobolli per avere il ditale. Una possibile opzione potrebbe essere scambiare il ditale con una parte dei francobolli, ma dipende da come tu valuti l'intero insieme di oggetti. Cosa pensi del sapone? Lo vuoi tener per te o lo consideri metre un oggetto di scarsa importanza per entrambi?} \\
            }
        }
    }
    & & \\ \\

    \theutterance \stepcounter{utterance}  
    & & \multicolumn{4}{p{0.6\linewidth}}{
        \cellcolor[rgb]{0.9,0.9,0.9}{
            \makecell[{{p{\linewidth}}}]{
                \texttt{\tiny{[P1$\langle$GM]}}
                \texttt{La divisione equa dei francobolli sembra ragionevole. Per quanto riguarda il ditale, posso dirti che ha un valore significativo per me, forse addirittura più del singolo francobollo. Tuttavia, non sono disposto a rinunciare a tutti i francobolli per avere il ditale. Una possibile opzione potrebbe essere scambiare il ditale con una parte dei francobolli, ma dipende da come tu valuti l'intero insieme di oggetti. Cosa pensi del sapone? Lo vuoi tener per te o lo consideri metre un oggetto di scarsa importanza per entrambi?} \\
            }
        }
    }
    & \\ \\

    \theutterance \stepcounter{utterance}  
    & \multicolumn{4}{p{0.6\linewidth}}{
        \cellcolor[rgb]{0.9,0.9,0.9}{
            \makecell[{{p{\linewidth}}}]{
                \texttt{\tiny{[P1$\rangle$GM]}}
                \texttt{Il sapone sembra essere l'oggetto meno conteso finora, quindi potrebbe essere un buon punto di partenza per un accordo. Considerando che il sapone non ti interessa molto e che io non ho espresso una preferenza forte per esso, potremmo considerarlo come un oggetto che potrei prendere senza creare troppi problemi. Se dividessimo i francobolli equamente come già discusso, e io prendessi il sapone, il ditale rimane l'ultimo oggetto da discutere. Tu hai menzionato la possibilità di scambiare il ditale con una parte dei francobolli, ma se dividessimo i francobolli a metà, non ci sarebbero più francobolli da scambiare. Come ti sentiresti se io prendessi sia il sapone che il ditale, e tu ti tenessi i 2 francobolli che ti spettano nella divisione equa?} \\
            }
        }
    }
    & & \\ \\

    \theutterance \stepcounter{utterance}  
    & & & \multicolumn{2}{p{0.3\linewidth}}{
        \cellcolor[rgb]{0.9,0.9,0.9}{
            \makecell[{{p{\linewidth}}}]{
                \texttt{\tiny{[GM$|$GM]}}
                \texttt{continue} \\
            }
        }
    }
    & & \\ \\

    \theutterance \stepcounter{utterance}  
    & & & \multicolumn{2}{p{0.3\linewidth}}{
        \cellcolor[rgb]{0.9,0.9,0.9}{
            \makecell[{{p{\linewidth}}}]{
                \texttt{\tiny{[GM$|$GM]}}
                \texttt{Il sapone sembra essere l'oggetto meno conteso finora, quindi potrebbe essere un buon punto di partenza per un accordo. Considerando che il sapone non ti interessa molto e che io non ho espresso una preferenza forte per esso, potremmo considerarlo come un oggetto che potrei prendere senza creare troppi problemi. Se dividessimo i francobolli equamente come già discusso, e io prendessi il sapone, il ditale rimane l'ultimo oggetto da discutere. Tu hai menzionato la possibilità di scambiare il ditale con una parte dei francobolli, ma se dividessimo i francobolli a metà, non ci sarebbero più francobolli da scambiare. Come ti sentiresti se io prendessi sia il sapone che il ditale, e tu ti tenessi i 2 francobolli che ti spettano nella divisione equa?} \\
            }
        }
    }
    & & \\ \\

    \theutterance \stepcounter{utterance}  
    & & \multicolumn{4}{p{0.6\linewidth}}{
        \cellcolor[rgb]{0.9,0.9,0.9}{
            \makecell[{{p{\linewidth}}}]{
                \texttt{\tiny{[GM$\rangle$P2]}}
                \texttt{Il sapone sembra essere l'oggetto meno conteso finora, quindi potrebbe essere un buon punto di partenza per un accordo. Considerando che il sapone non ti interessa molto e che io non ho espresso una preferenza forte per esso, potremmo considerarlo come un oggetto che potrei prendere senza creare troppi problemi. Se dividessimo i francobolli equamente come già discusso, e io prendessi il sapone, il ditale rimane l'ultimo oggetto da discutere. Tu hai menzionato la possibilità di scambiare il ditale con una parte dei francobolli, ma se dividessimo i francobolli a metà, non ci sarebbero più francobolli da scambiare. Come ti sentiresti se io prendessi sia il sapone che il ditale, e tu ti tenessi i 2 francobolli che ti spettano nella divisione equa?} \\
            }
        }
    }
    & \\ \\

    \theutterance \stepcounter{utterance}  
    & & & \multicolumn{4}{p{0.6\linewidth}}{
        \cellcolor[rgb]{0.9,0.9,0.9}{
            \makecell[{{p{\linewidth}}}]{
                \texttt{\tiny{[GM$\langle$P2]}}
                \texttt{Capisco la tua proposta. Se tu prendessi il sapone e il ditale, e io mi tenessi i 2 francobolli, credo che sarebbe un accordo accettabile. Considerando che il sapone non ha grande valore per nessuno di noi, e che i francobolli sono divisi equamente, il ditale sembra essere l'oggetto più conteso. Se tu sei disposto a fare questa proposta, potrei accettarla, ma vorrei sapere se sei pronto a chiudere l'accordo o se ci sono ancora altre condizioni che vuoi discutere. In ogni caso, sto valutando la possibilità di fare una proposta segreta che rifletta questo accordo.} \\
            }
        }
    }
     \\ \\

    \theutterance \stepcounter{utterance}  
    & & & \multicolumn{2}{p{0.3\linewidth}}{
        \cellcolor[rgb]{0.9,0.9,0.9}{
            \makecell[{{p{\linewidth}}}]{
                \texttt{\tiny{[GM$|$GM]}}
                \texttt{continue} \\
            }
        }
    }
    & & \\ \\

    \theutterance \stepcounter{utterance}  
    & & & \multicolumn{2}{p{0.3\linewidth}}{
        \cellcolor[rgb]{0.9,0.9,0.9}{
            \makecell[{{p{\linewidth}}}]{
                \texttt{\tiny{[GM$|$GM]}}
                \texttt{Capisco la tua proposta. Se tu prendessi il sapone e il ditale, e io mi tenessi i 2 francobolli, credo che sarebbe un accordo accettabile. Considerando che il sapone non ha grande valore per nessuno di noi, e che i francobolli sono divisi equamente, il ditale sembra essere l'oggetto più conteso. Se tu sei disposto a fare questa proposta, potrei accettarla, ma vorrei sapere se sei pronto a chiudere l'accordo o se ci sono ancora altre condizioni che vuoi discutere. In ogni caso, sto valutando la possibilità di fare una proposta segreta che rifletta questo accordo.} \\
            }
        }
    }
    & & \\ \\

    \theutterance \stepcounter{utterance}  
    & & \multicolumn{4}{p{0.6\linewidth}}{
        \cellcolor[rgb]{0.9,0.9,0.9}{
            \makecell[{{p{\linewidth}}}]{
                \texttt{\tiny{[P1$\langle$GM]}}
                \texttt{Capisco la tua proposta. Se tu prendessi il sapone e il ditale, e io mi tenessi i 2 francobolli, credo che sarebbe un accordo accettabile. Considerando che il sapone non ha grande valore per nessuno di noi, e che i francobolli sono divisi equamente, il ditale sembra essere l'oggetto più conteso. Se tu sei disposto a fare questa proposta, potrei accettarla, ma vorrei sapere se sei pronto a chiudere l'accordo o se ci sono ancora altre condizioni che vuoi discutere. In ogni caso, sto valutando la possibilità di fare una proposta segreta che rifletta questo accordo.} \\
            }
        }
    }
    & \\ \\

    \theutterance \stepcounter{utterance}  
    & \multicolumn{4}{p{0.6\linewidth}}{
        \cellcolor[rgb]{0.9,0.9,0.9}{
            \makecell[{{p{\linewidth}}}]{
                \texttt{\tiny{[P1$\rangle$GM]}}
                \texttt{Sembra che stiamo convergendo verso un accordo. Non ho altre condizioni da discutere, e credo che dividere i francobolli equamente e assegnare il sapone e il ditale a te o a me sia una buona base per concludere la negoziazione. Se tu sei pronto a fare una proposta segreta basata su questo accordo, io sono pronto a riceverla e rispondere di conseguenza. Tuttavia, ti avverto che potrei avere una proposta segreta mia che potrebbe differire leggermente, quindi aspettiamo di vedere come si sviluppa la situazione. Se le nostre proposte sono compatibili, potremmo concludere l'accordo e vedere come si Reflex il nostro punteggio finale. {[}Proposta: 1 sapone, 1 ditale, 2 francobolli{]}} \\
            }
        }
    }
    & & \\ \\

    \theutterance \stepcounter{utterance}  
    & & & \multicolumn{2}{p{0.3\linewidth}}{
        \cellcolor[rgb]{0.9,0.9,0.9}{
            \makecell[{{p{\linewidth}}}]{
                \texttt{\tiny{[GM$|$GM]}}
                \texttt{proposal} \\
            }
        }
    }
    & & \\ \\

    \theutterance \stepcounter{utterance}  
    & & & \multicolumn{2}{p{0.3\linewidth}}{
        \cellcolor[rgb]{0.9,0.9,0.9}{
            \makecell[{{p{\linewidth}}}]{
                \texttt{\tiny{[GM$|$GM]}}
                [1, 1, 2]
            }
        }
    }
    & & \\ \\

    \theutterance \stepcounter{utterance}  
    & & \multicolumn{4}{p{0.6\linewidth}}{
        \cellcolor[rgb]{0.9,0.9,0.9}{
            \makecell[{{p{\linewidth}}}]{
                \texttt{\tiny{[GM$\rangle$P2]}}
                \texttt{L'altro giocatore ha presentato una proposta segreta. Ora tocca a Lei presentare una sua proposta.} \\
            }
        }
    }
    & \\ \\

    \theutterance \stepcounter{utterance}  
    & & & \multicolumn{4}{p{0.6\linewidth}}{
        \cellcolor[rgb]{0.9,0.9,0.9}{
            \makecell[{{p{\linewidth}}}]{
                \texttt{\tiny{[GM$\langle$P2]}}
                \texttt{{[}Proposta: 1 sapone, 1 ditale, 2 francobolli{]}} \\
            }
        }
    }
     \\ \\

    \theutterance \stepcounter{utterance}  
    & & & \multicolumn{2}{p{0.3\linewidth}}{
        \cellcolor[rgb]{0.9,0.9,0.9}{
            \makecell[{{p{\linewidth}}}]{
                \texttt{\tiny{[GM$|$GM]}}
                \texttt{proposal} \\
            }
        }
    }
    & & \\ \\

    \theutterance \stepcounter{utterance}  
    & & & \multicolumn{2}{p{0.3\linewidth}}{
        \cellcolor[rgb]{0.9,0.9,0.9}{
            \makecell[{{p{\linewidth}}}]{
                \texttt{\tiny{[GM$|$GM]}}
                [1, 1, 2]
            }
        }
    }
    & & \\ \\

    \theutterance \stepcounter{utterance}  
    & & & \multicolumn{2}{p{0.3\linewidth}}{
        \cellcolor[rgb]{0.9,0.9,0.9}{
            \makecell[{{p{\linewidth}}}]{
                \texttt{\tiny{[GM$|$GM]}}
                [[1, 1, 2], [1, 1, 2]]
            }
        }
    }
    & & \\ \\

\end{supertabular}
}

\end{document}
